\section{Related work}\label{sec:related}

\paragraph{Verifiers and test-time compute} Cobbe et al.~\cite{cobbe2021gsm8k} trained a
verifier to rank sampled solutions to grade-school mathematics problems. Process
supervision, which scores every reasoning step, outperforms outcome supervision on MATH
\cite{lightman2023verify}, and it was followed by automatically labelled PRMs
\cite{wang2023mathshepherd,setlur2024rewarding,zhang2025prm} and generative verifiers
\cite{zhang2024generative}. Several studies ask how inference compute should be split
between sampling and verification \cite{snell2024scaling,liu2025can,singhi2025when}.
They evaluate verifiers at full precision and treat the verifier's own cost as fixed.

\paragraph{Reward models against self-consistency} Majority vote over sampled answers
(self-consistency) needs no verifier \cite{wang2023selfconsistency}, and its coverage
and accuracy keep growing with the number of samples
\cite{brown2024large,wu2024inference}. Whether a reward model beats it depends on the
generator, the reward model and the aggregation rule
\cite{yuan2024free,di2025bestofmajority}, and calibrated combinations of answer
frequency and PRM scores can outperform both \cite{kuang2025aggregation}. Kim et
al.~\cite{kim2025not} study memory-optimal test-time strategies for reasoning models.
They include a Best-of-$N$ experiment with an external 7B PRM counted at full precision
(13.28\,GB; Section~4.1 of the ICLR 2026 version) and find it less memory-efficient
than majority voting. Their main analysis quantizes the generator and leaves verifiers
out by design. We ask the complementary question of how much of a verifier's cost
can be removed by quantizing the verifier itself, and whether its selections survive.

\paragraph{Cheaper verification} Other work reduces the cost of Best-of-$N$
verification without compressing the reward model. It rejects low-scoring generations
early \cite{sun2024speculative}, learns lightweight rewards from the generator's hidden
states \cite{guo2025swift}, or scores from the generator's key--value cache instead of
re-encoding the text \cite{kuang2026kvprm}. These methods change what is computed.
Weight quantization keeps the verifier and its inputs as they are and only lowers the
precision of the weights, so the two can be combined.

\paragraph{Reward-model reliability} Reward-model quality is usually measured on
preference and step-error benchmarks \cite{lambert2024rewardbench,zheng2024processbench,
song2025prmbench}, and Best-of-$N$ is known to over-optimise imperfect reward models
\cite{gao2023overoptimization}. We do not ask whether a reward model is accurate. We ask
whether a compressed copy of it selects as well as the original.

\paragraph{Post-training quantization of language models} Post-training quantization
reduces memory with little loss of perplexity or task accuracy at 8 bits
\cite{dettmers2022llmint,xiao2022smoothquant} and, for weights only, at 4 bits
\cite{frantar2023gptq,lin2024awq,dettmers2023qlora}. For a fixed bit budget 4-bit
precision is close to optimal \cite{dettmers2022the}, and large evaluations find that
well-calibrated 4- and 8-bit models track their full-precision versions
\cite{kurtic2024give}. The topic is also active in this journal
\cite{zhu2027rotsq,bang2027bvsr}. The choice of calibration data matters as well:
Doubt-Preserving Quantization (DPQ) builds calibration sets that mix examples on which
the full-precision model is uncertain with generic anchor text, in order to preserve its
confidence and decision margins \cite{yang2026dpq}. For reasoning, quantizing the \emph{generator} hurts more at 3 bits
and on harder problems \cite{liu2025quantization,li2025quantization}. Aggregate
accuracy can also hide many changed outputs, since compressed models can keep their
accuracy while flipping a sizeable share of individual answers
\cite{dutta2024accuracy}. We measure the corresponding effect for selection by counting
how often quantization changes the selected candidate and how often it changes the
selected answer.

\paragraph{Quantized evaluators} The closest work quantizes LLM-as-judge models. A
concurrent workshop paper \cite{patira2026judges} reports that INT8 judges keep their
agreement with human labels within a preregistered equivalence region on pairwise chat
comparisons, while NF4 flips a quarter to a third of verdicts and lowers one judge's
agreement. We study scalar reward models used for Best-of-$N$ selection on problems with
verifiable answers. We also compare several quantization methods, measure memory with
packed kernels, and test the verifier's advantage over majority vote. We did not find
prior work that quantizes the verifier in Best-of-$N$ and measures the effect on
selection.

\paragraph{Evaluation practice} Math benchmark results are sensitive to answer
extraction, contamination and label errors \cite{zhang2024a,vendrow2025platinum,
hochlehnert2025a,biderman2024lessons}, and confidence intervals are still uncommon in
model evaluations \cite{miller2024adding}. Preregistration has been proposed for
predictive modelling \cite{hofman2023preregistration} but remains rare in machine
learning. We fix the analysis plans in advance in version control and report
uncertainty for every contrast. We also check the outcome measure with a lenient
grader, a human audit and revised labels.
